% NOTACION: en todo el informe la incertidumbre típica se denota sigma(.),
% por continuidad con el marco teórico del capítulo 1. Equivale a u(.) en la
% notación de la GUM (JCGM 100:2008). Si más adelante se prefiere la notación
% GUM en la Tabla 6, el generador está en la celda del notebook que exporta
% t6.tex, en las líneas marcadas [GUM].

\paragraph{Estimación de las componentes de incertidumbre del pie de rey \texorpdfstring{$\sigma(L_i)$}{sigma(Li)}}
Para cuantificar las fuentes de variabilidad que afectan directamente la lectura del pie de rey, se evalúan las componentes geométricas, estadísticas y de resolución conforme al modelo de varianza de la ecuación~\eqref{eq:discrete-std-x-i}. Para cada componente acotada se aplica la distribución rectangular deducida en el marco teórico (ecuación~\eqref{eq:sigma-rectangular}), donde $\sigma = a/\sqrt{3}$ con $a$ la semiamplitud del intervalo.

La componente asociada al error de contacto de la superficie parcial, $\sigma(L_{i_E})$, se dedujo en la Sección~\ref{sec:pie-rey-reduccion} (ecuación~\eqref{eq:sigma-LiE}). A continuación se estiman las restantes.

\paragraph{Error de repetibilidad de contacto de la superficie parcial ($R$)}
Para cuantificar la dispersión aleatoria experimental (evaluación Tipo A), se seleccionan $n = 5$ mediciones repetidas en el punto de mayor error nominal. De acuerdo con las ecuaciones~\eqref{eq:std-s-eq} y~\eqref{eq:std-relative}, la desviación del promedio de $n$ observaciones independientes se reduce por el factor $1/\sqrt{n}$:

\begin{equation}
    \sigma(L_{i_R}) = \frac{s_R}{\sqrt{n}}
    \label{eq:sigma-LiR}
\end{equation}

donde la desviación estándar experimental de la muestra $s_R$ corresponde a:

\begin{equation}
    s_R = \sqrt{\frac{1}{n-1} \sum_{p=1}^n (R_p - \bar{R})^2}
    \label{eq:desv-estandar-muestral}
\end{equation}

siendo $R_p$ cada lectura individual y $\bar{R}$ la media aritmética muestral.

Las cinco repeticiones se efectuaron en el punto de $250$ mm, seleccionado por ser
el de mayor lectura de la serie (ante el empate en $E$ entre $100$ mm y $200$ mm se
optó por el extremo del alcance, donde el error de repetibilidad suele ser mayor):

\begin{table}[H]
\centering
\caption{Repetibilidad en el punto de mayor error de contacto ($250$ mm), $n = 5$ mediciones repetidas.}
\label{tab:pie-rey-repetibilidad}
\begin{tabular}{rr}
\toprule
$p$ & $X_p$ (mm) \\
\midrule
1 & 250.10 \\
2 & 250.05 \\
3 & 250.10 \\
4 & 250.10 \\
5 & 250.05 \\
\bottomrule
\end{tabular}
\end{table}


El rango de las cinco lecturas, $R = \max R_p - \min R_p = 50\ \mu\text{m}$, y la
desviación estándar experimental $s_R = 27.39\ \mu\text{m}$ conducen a

\begin{equation}
    \therefore\ \sigma(L_{i_R}) = \frac{27.39}{\sqrt{5}}\ \mu\text{m}
    \approx 12.25\ \mu\text{m}
    \label{eq:sigma-LiR-valor}
\end{equation}

Conviene notar que las cinco repeticiones fueron tomadas en una tanda distinta de
las mediciones de la Tabla~\ref{tab:pie-rey-datos}, donde las tres lecturas de
$250$ mm fueron $250.00$ mm. La diferencia de $0.08$ mm entre la media de las
repeticiones ($\bar{R} = 250.08$ mm) y aquella lectura refleja una deriva
instrumental entre ambas series, no dispersión aleatoria, y por ello no se suma a
$\sigma(L_i)$: la repetibilidad cuantifica únicamente la dispersión interna de la
tanda repetida, tal como define la ecuación~\eqref{eq:sigma-LiR}.

\paragraph{Error por cambio de escala ($S$)}
Representa la discrepancia al alternar entre la escala de medición de exteriores e interiores ($S_{E-I}$) o exteriores y profundidad ($S_{E-P}$). Tomando el caso más desfavorable $S = \max(S_{E-I}, S_{E-P})$ bajo el modelo de la ecuación~\eqref{eq:sigma-rectangular}:

\begin{equation}
    \sigma(L_{i_S}) = \frac{\max(S_{E-I}, S_{E-P})}{2\sqrt{3}}
    \label{eq:sigma-LiS}
\end{equation}

El ensayo se realizó sobre un cilindro patrón de $10$ mm, midiéndolo en tres
posiciones sobre la generatriz (superior, central e inferior) con cada una de las
tres aplicaciones del instrumento: exteriores, interiores y profundidad. La
Tabla~\ref{tab:pie-rey-cilindro} recoge las lecturas y los errores de indicación
correspondientes.

\begin{table}[H]
\centering
\caption{Cambio de escala sobre el cilindro patrón de $10$ mm. Las columnas corresponden a la medición de exteriores, interiores y profundidad, respectivamente; las filas, a la posición sobre cada superficie.}
\label{tab:pie-rey-cilindro}
\begin{tabular}{lrrrrrr}
\toprule
Posición $k$ & $X_{k,1}$ (mm) & $X_{k,2}$ (mm) & $X_{k,3}$ (mm) & $D_{k,1}$ ($\mu$m) & $D_{k,2}$ ($\mu$m) & $D_{k,3}$ ($\mu$m) \\
\midrule
1 (superior) & 10.00 & 10.00 & 10.05 & 0 & 0 & 50 \\
2 (central) & 10.00 & 10.05 & 10.05 & 0 & 50 & 50 \\
3 (inferior) & 10.00 & 10.00 & 10.00 & 0 & 0 & 0 \\
\bottomrule
\end{tabular}
\end{table}


Promediando las diferencias respecto de la medición de exteriores se obtiene
$S_{E-I} = 16.7\ \mu\text{m}$ y $S_{E-P} = 33.3\ \mu\text{m}$. El caso más
desfavorable es la medición de profundidad:

\begin{equation}
    \therefore\ \sigma(L_{i_S}) =
    \frac{33.3}{2\sqrt{3}}\ \mu\text{m} \approx 9.62\ \mu\text{m}
    \label{eq:sigma-LiS-valor}
\end{equation}

\paragraph{Errores geométricos de contacto y alineamiento ($L$, $J$, $K$)}
Siguiendo análogamente el modelo de la ecuación~\eqref{eq:sigma-rectangular}, las incertidumbres estándar por desviaciones de contacto y cruce resultan:

\begin{align}
    \sigma(L_{i_L}) &= \frac{L}{2\sqrt{3}} \label{eq:sigma-LiL} \\
    \sigma(L_{i_J}) &= \frac{J}{2\sqrt{3}} \label{eq:sigma-LiJ} \\
    \sigma(L_{i_K}) &= \frac{K}{2\sqrt{3}} \label{eq:sigma-LiK}
\end{align}

La medición del cilindro sobre la superficie de contacto lineal (aristas de las
mordazas sobre una generatriz) arrojó $L = 10\ \mu\text{m}$; sobre superficie
completa (bloques), $J = 0$; y en el cruce de las superficies de medición de
interiores, $K = 0$. Por tanto

\begin{equation}
    \sigma(L_{i_L}) = \frac{10}{2\sqrt{3}}\ \mu\text{m} \approx 2.89\ \mu\text{m},
    \qquad \sigma(L_{i_J}) = \sigma(L_{i_K}) = 0
    \label{eq:sigma-LiLJK-valor}
\end{equation}

\paragraph{Incertidumbre por resolución del pie de rey ($r$)}
El límite de discriminación instrumental define un intervalo de semiamplitud $\pm r/2$. Para escalas vernier se toma $r = d/m$ (con $d = 0.05\text{ mm}$, la división de escala del 530-115, y $m = 2$), obteniéndose mediante la ecuación~\eqref{eq:sigma-rectangular}:

\begin{equation}
    \sigma(L_{i_r}) = \frac{r}{\sqrt{3}} = \frac{d / m}{\sqrt{3}}
    \label{eq:sigma-Lir}
\end{equation}

Sustituyendo los valores del instrumento:

\begin{equation}
    \therefore\ \sigma(L_{i_r}) =
    \frac{0.05/2}{\sqrt{3}}\ \text{mm} \approx 14.43\ \mu\text{m}
    \label{eq:sigma-Lir-valor}
\end{equation}

\paragraph{Varianza combinada de la indicación $\sigma^2(L_i)$}
Asumiendo independencia estadística entre los factores evaluados ($\text{Cov} = 0$, conforme a la ecuación~\eqref{eq:error-std-values}), la varianza total de la indicación del pie de rey resulta de la suma en cuadratura de las ecuaciones~\eqref{eq:sigma-LiE} a~\eqref{eq:sigma-Lir}:

\begin{equation}
    \sigma^2(L_i) = \sigma^2(L_{i_E}) + \sigma^2(L_{i_R}) + \sigma^2(L_{i_S}) + \sigma^2(L_{i_L}) + \sigma^2(L_{i_J}) + \sigma^2(L_{i_K}) + \sigma^2(L_{i_r})
    \label{eq:sigma-Li-total}
\end{equation}

\subsubsection{Incertidumbre debida a los patrones de referencia \texorpdfstring{$\sigma(L_{BP})$}{sigma(LBP)}}

La longitud de referencia $L_{BP}$ incorpora dos contribuciones independientes:

\paragraph{Incertidumbre de calibración y combinación de bloques}
A partir de la incertidumbre expandida $U_{cal}$ reportada en el certificado para un factor de cobertura $k = 2$, la varianza combinada para una serie de $M$ bloques patrón acoplados es:

\begin{equation}
    \sigma^2(L_{BP_{cal}}) = \sum_{q=1}^M \left(\frac{U_{Lq,cal}}{k}\right)^2
    \label{eq:sigma-BP-cal}
\end{equation}

\paragraph{Deriva temporal de los bloques patrón}
Modelando la variación máxima permisible anual $U_{der}$ estipulada por la norma ISO 3650 mediante la distribución rectangular de la ecuación~\eqref{eq:sigma-rectangular}:

\begin{equation}
    \sigma^2(L_{BP_{der}}) = \sum_{q=1}^M \left(\frac{U_{Lq,der}}{\sqrt{3}}\right)^2
    \label{eq:sigma-BP-der}
\end{equation}

\paragraph{Varianza total del bloque patrón $\sigma^2(L_{BP})$}
Combinando ambas componentes independientes:

\begin{equation}
    \sigma^2(L_{BP}) = \sigma^2(L_{BP_{cal}}) + \sigma^2(L_{BP_{der}})
    \label{eq:sigma-LBP-total}
\end{equation}

El certificado de los bloques patrón no reporta una incertidumbre individual por
pieza, sino la expresión funcional con que el laboratorio la declara en función de
la longitud nominal $l$:

\begin{equation}
    U_{cal}(l) = 2\sqrt{a^2 + b^2 l^2},
    \qquad a = 41.46\ \text{nm},\quad b = 0.55\ \text{nm/mm},
    \label{eq:Ucal-bloque}
\end{equation}

con factor de cobertura $k = 2$. La deriva se estima con la variación máxima
permisible anual de la norma ISO~3650 para bloques de grado~1,
$U_{der}(l) = \pm(0.05\ \mu\text{m} + 0.5\times10^{-6}l)$. Evaluando ambas
expresiones sobre los cuatro bloques de la pila de $250$ mm ($25$, $50$, $75$ y
$100$ mm) resulta

\begin{equation}
    \sigma(L_{BP_{cal}}) = 0.1120\ \mu\text{m},
    \qquad \sigma(L_{BP_{der}}) = 0.0578\ \mu\text{m},
    \label{eq:sigma-BP-valor}
\end{equation}

\begin{equation}
    \therefore\ \sigma(L_{BP}) = \sqrt{0.1120^2 + 0.0578^2}\ \mu\text{m}
    \approx 0.126\ \mu\text{m}
    \label{eq:sigma-LBP-total-valor}
\end{equation}

\subsubsection{Incertidumbres asociadas a los parámetros térmicos}

Para las condiciones ambientales de ensayo, la temperatura registrada con el
termómetro de contacto fue de $21.5\ ^\circ\text{C}$ en los cinco puntos de
calibración (Tabla~\ref{tab:pie-rey-datos}), sin variación apreciable entre ellos. De
ahí resulta una desviación respecto de la temperatura de referencia de
$\Delta T_i = 21.5 - 20 = 1.5\ ^\circ\text{C}$, y una diferencia nula entre el bloque
patrón y el pie de rey, $\delta T = 0$, ya que ambos permanecieron en el mismo
ambiente durante la medición.

\paragraph{Coeficientes de dilatación térmica $\sigma^2(\alpha_{BP})$ y $\sigma^2(\delta\alpha)$}
Para bloques y pie de rey de acero con semiamplitudes de tolerancia $\Delta\alpha_{BP}$ y $\Delta\alpha_i$, la aplicación de la ecuación~\eqref{eq:sigma-rectangular} proporciona:

\begin{align}
    \sigma^2(\alpha_{BP}) &= \frac{\Delta\alpha_{BP}^2}{3}, \label{eq:sigma-alpha-BP} \\
    \sigma^2(\delta\alpha) &= \frac{\Delta\alpha_{BP}^2}{3} + \frac{\Delta\alpha_i^2}{3} \label{eq:sigma-dalpha}
\end{align}

El acero presenta $\alpha = 11.5\times10^{-6}\ ^\circ\text{C}^{-1}$ con una dispersión
de $\Delta\alpha = 1\times10^{-6}\ ^\circ\text{C}^{-1}$, igual para el bloque y para el
pie de rey por ser ambos del mismo material. Sustituyendo en las
ecuaciones~\eqref{eq:sigma-alpha-BP} y~\eqref{eq:sigma-dalpha}:

\begin{equation}
    \sigma(\alpha_{BP}) = \frac{1\times10^{-6}}{\sqrt{3}}\ ^\circ\text{C}^{-1}
    \approx 5.77\times10^{-7}\ ^\circ\text{C}^{-1}
    \label{eq:sigma-alpha-BP-valor}
\end{equation}

\begin{equation}
    \sigma(\delta\alpha) = \frac{\sqrt{2}\times10^{-6}}{\sqrt{3}}\ ^\circ\text{C}^{-1}
    \approx 8.16\times10^{-7}\ ^\circ\text{C}^{-1}
    \label{eq:sigma-dalpha-valor}
\end{equation}

\paragraph{Desviaciones y gradientes de temperatura $\sigma^2(\Delta T_i)$ y $\sigma^2(\delta T)$}
Para la desviación respecto de la temperatura de referencia intervienen la
incertidumbre del certificado del termómetro ($UT_{cert}/k$), la deriva instrumental
($UT_{der}/\sqrt{3}$), su resolución ($d_T / (m_T \sqrt{3})$) y la fluctuación
observada durante el ensayo ($\Delta t_{A_i}/\sqrt{3}$):

\begin{equation}
    \sigma^2(\Delta T_i) = \left(\frac{UT_{cert}}{k}\right)^2 + \frac{\Delta t_{A_i}^2}{3} + \left(\frac{d_T / m_T}{\sqrt{3}}\right)^2 + \frac{UT_{der}^2}{3}
    \label{eq:sigma-deltaTi}
\end{equation}

El termómetro de contacto del laboratorio tiene una exactitud de
$\pm 0.5\ ^\circ\text{C}$, una deriva de $0.5\ ^\circ\text{C}$ y una división de
escala $d_T = 0.2\ ^\circ\text{C}$ subdividida en $m_T = 5$ partes. Como la
temperatura registrada fue la misma en los cinco puntos, la fluctuación observada es
nula ($\Delta t_{A_i} = 0$). Sustituyendo en la ecuación~\eqref{eq:sigma-deltaTi}:

\begin{equation}
    \sigma(\Delta T_i) = \sqrt{\left(\frac{0.5}{2}\right)^2 + 0
    + \left(\frac{0.2/5}{\sqrt{3}}\right)^2 + \frac{0.5^2}{3}}\ ^\circ\text{C}
    \approx 0.383\ ^\circ\text{C}
    \label{eq:sigma-deltaTi-valor}
\end{equation}

De manera análoga, para el gradiente térmico diferencial entre los bloques y el pie de rey ($\delta T$):

\begin{equation}
    \sigma^2(\delta T) = \left(\frac{UT_{cert}}{k}\right)^2 + \frac{\Delta t^2}{3} + \left(\frac{d_T / m_T}{\sqrt{3}}\right)^2 + \frac{UT_{der}^2}{3}
    \label{eq:sigma-deltaT}
\end{equation}

donde $\Delta t = \max(\Delta t_1, \Delta t_2)$ representa la máxima diferencia térmica entre ambos cuerpos. Al haberse mantenido bloque y pie de rey en el mismo ambiente, esa diferencia también es nula, y por tanto

\begin{equation}
    \sigma(\delta T) \approx 0.383\ ^\circ\text{C}
    \label{eq:sigma-deltaT-valor}
\end{equation}

\paragraph{Contribución efectiva de los términos térmicos}
Las incertidumbres anteriores no se transmiten al resultado, porque los coeficientes
de sensibilidad que las multiplican en la ecuación~\eqref{eq:sigma-final} se anulan:
ambos cuerpos son de acero, de modo que $\delta\alpha = \alpha_i - \alpha_{BP} = 0$, y
ambos están a la misma temperatura, de modo que $\delta T = 0$. El único término
térmico que sobrevive es $[L_i\,\Delta T_i]^2 \sigma^2(\delta\alpha)$, que evaluado en
el extremo del alcance ($L_i = 250$ mm, $\Delta T_i = 1.5\ ^\circ\text{C}$) vale

\begin{equation}
    (250 \times 1.5)^2 \left(8.16\times10^{-7}\right)^2\ \mu\text{m}^2
    \approx 9.4\times10^{-8}\ \mu\text{m}^2,
    \label{eq:termino-termico}
\end{equation}

es decir diez órdenes de magnitud por debajo de $\sigma^2(L_i) = 667.6\ \mu\text{m}^2$.
Los parámetros térmicos son, en consecuencia, irrelevantes para la incertidumbre del
error de indicación de este instrumento en las condiciones de ensayo.

\subsubsection{Incertidumbre estándar combinada y expandida final}

Sustituyendo las componentes~\eqref{eq:sigma-Li-total},~\eqref{eq:sigma-LBP-total},~\eqref{eq:sigma-dalpha},~\eqref{eq:sigma-deltaTi} y~\eqref{eq:sigma-deltaT} en la ley de propagación de varianzas de la ecuación~\eqref{eq:sigma-final}, la varianza combinada del error de indicación queda expresada en forma completamente desarrollada como:

\begin{align}
    \sigma_c^2(E) &=
    [1 + \delta\alpha \, \Delta T_i]^2
    \Bigg[
        \left(\frac{E}{2\sqrt{3}}\right)^2
        + \frac{s_R^2}{n}
        + \left(\frac{\max(S_{E-I}, S_{E-P})}{2\sqrt{3}}\right)^2 \nonumber \\
    &\qquad\qquad
        + \left(\frac{L}{2\sqrt{3}}\right)^2
        + \left(\frac{J}{2\sqrt{3}}\right)^2
        + \left(\frac{K}{2\sqrt{3}}\right)^2
        + \left(\frac{d/m}{\sqrt{3}}\right)^2
    \Bigg] \nonumber \\[4pt]
    &\quad + [-1 + \alpha_{BP} \, \delta T]^2
    \Bigg[
        \sum_{q=1}^M \left(\frac{U_{Lq,cal}}{k}\right)^2
        + \sum_{q=1}^M \left(\frac{U_{Lq,der}}{\sqrt{3}}\right)^2
    \Bigg] \nonumber \\[4pt]
    &\quad + [L_i \, \Delta T_i]^2
    \Bigg[
        \frac{\Delta\alpha_{BP}^2}{3} + \frac{\Delta\alpha_i^2}{3}
    \Bigg] \nonumber \\[4pt]
    &\quad + [L_i \, \delta\alpha]^2
    \Bigg[
        \left(\frac{UT_{cert}}{k}\right)^2
        + \frac{\Delta t_{A_i}^2}{3}
        + \left(\frac{d_T / m_T}{\sqrt{3}}\right)^2
        + \frac{UT_{der}^2}{3}
    \Bigg] \nonumber \\[4pt]
    &\quad + [L_{BP} \, \delta T]^2
    \Bigg[
        \frac{\Delta\alpha_{BP}^2}{3}
    \Bigg] \nonumber \\[4pt]
    &\quad + [L_{BP} \, \alpha_{BP}]^2
    \Bigg[
        \left(\frac{UT_{cert}}{k}\right)^2
        + \frac{\Delta t^2}{3}
        + \left(\frac{d_T / m_T}{\sqrt{3}}\right)^2
        + \frac{UT_{der}^2}{3}
    \Bigg] \label{eq:sigma-cuad-expandida}
\end{align}

Finalmente, la incertidumbre expandida $U(E)$ para el error de indicación del pie de rey se reporta mediante:

\begin{equation}
    U(E) = k \cdot \sigma_c(E) = k \sqrt{\sigma_c^2(E)}
    \label{eq:incertidumbre-expandida}
\end{equation}

adoptando un factor de cobertura $k = 2$, correspondiente a un nivel de confianza bilateral aproximado del 95\,\%.
